% !TeX root = main.tex

\section{Inleiding}

In de hedendaagse topsport wordt het verschil tussen winst en verlies vaak 
bepaald door fracties van seconden en enkele centimeters. Waar tijdwaarneming 
zich vroeger beperkte tot het vastleggen van de finish, vereisen moderne 
analyses en live tracking een continue en accurate positiebepaling van atleten 
en voertuigen. \\

MYLAPS Sports Timing speelt hierin wereldwijd een belangrijke rol. Het van 
oorsprong Nederlandse bedrijf levert al decennialang de hardware en software 
voor duizenden evenementen, variërend van lokale hardloopwedstrijden tot de 
Olympische Spelen en grote motorsportkampioenschappen. De kern van MYLAPS 
breidt zich echter steeds verder uit. Naast het accuraat registreren van 
tijden, ontstaat er vanuit de markt een sterke behoefte aan robuuste, continue 
real-time analyse van deelnemers gedurende een gehele race. \\

In opdracht van MYLAPS wordt in dit afstudeerproject gewerkt aan de volgende 
generatie van deze volgsystemen, de ProRTK tracker. Deze tracker heeft primair als doel ingezet te worden bij paardenraces, maar het systeem biedt de mogelijkheid om verder gebruijkt te worden bij andere disciplines, waaronder autosport. \\

De tracker maakt gebruik van `Real-Time Kinematic Global Navigation Satellite 
Systems' (RTK GNSS), waarbij de locatie tot op de centimeter nauwkeurig kan 
worden bepaald. In de praktijk ontbreekt echter de garantie op een constant 
betrouwbaar signaal. Signaaluitval en multipath effecten door tribunes en gebouwen kunnen leiden tot onbetrouwbare of ontbrekende positiedata. \\

De ProRTK tracker beschikt naast de RTK GNSS ontvanger ook over een `Inertial Measurement Unit' (IMU), bestaande uit een accelerometer en een gyroscoop. Deze IMU 
wordt in het huidige systeem echter niet ingezet voor de positiebepaling. Hierom is er vanuit MYLAPS de wens om te achterhalen of deze IMU (en eventuele aanvullende sensoren), doormiddel van sensorfusion (combinatie van meerdere sensordatastromen), meerwaarde kan bieden aan de huidige positiebepaling. De centrale vraag van dit onderzoek is of de aanwezige IMU ingezet kan worden om de robuustheid van de positiebepaling te verbeteren, en zo ja, in welke mate. \\

\todo[inline]{Opmaak van de hoofdstukken vermelden wanneer het verslag klaar is.}

\newpage